% Generated draft: requires mathematical and source review.
\documentclass[11pt]{article}
\usepackage[margin=25mm]{geometry}
\usepackage{fontspec}
\setmainfont{Latin Modern Roman}
\usepackage{amsmath,amssymb,mathtools,unicode-math}
\setmathfont{Latin Modern Math}
\usepackage{enumitem,needspace,xurl,hyperref}
\setlength{\parindent}{0pt}
\setlength{\parskip}{5pt}
\setlength{\emergencystretch}{3em}
\newcommand{\N}{\mathbb{N}}
\newcommand{\Z}{\mathbb{Z}}
\newcommand{\Q}{\mathbb{Q}}
\newcommand{\R}{\mathbb{R}}
\newcommand{\C}{\mathbb{C}}
\newcommand{\F}{\mathbb{F}}
\newcommand{\A}{\mathbb{A}}
\newcommand{\PP}{\mathbb{P}}
\newcommand{\E}{\mathbb{E}}
\newcommand{\eps}{\varepsilon}
\newcommand{\dd}{\,\mathrm{d}}
\begin{document}
\section*{UnsolvedMath ID: 30001883}
\textbf{Source problem code:} catalogue entry
\subsection{Definitions and mathematical statement}
Under the usual convention, a convex body in \(\R^3\) is a compact convex subset with nonempty interior; write \(\mathcal K(\R^3)\) for this class. For subsets \(A,D\subseteq\R^3\), their Minkowski sum is
\[
A+D=\{a+d:a\in A,\ d\in D\}.
\]
For \(c\in\R^3\) and \(r>0\), define the closed Euclidean ball
\[
B(c,r)=\{x\in\R^3:\|x-c\|_2\le r\}.
\]
The supplied problem asks whether adding a ball to any circle-free convex body in \(\R^3\) preserves circle-freeness. However, the supplied record does not define \emph{circle-free}. Let \(\operatorname{CF}(K)\) serve only as a placeholder for that intended, presently unspecified property. Under the closed, positive-radius ball convention above, the proposed assertion has the form
\[
\forall K\in\mathcal K(\R^3)\ \forall c\in\R^3\ \forall r>0,\qquad
\operatorname{CF}(K)\Longrightarrow\operatorname{CF}\bigl(K+B(c,r)\bigr).
\]
This is a formulation template, not a fully specified mathematical assertion: the intended definition of circle-freeness must be recovered from the cited material. The source record also does not explicitly specify whether balls are closed or open, or whether radius zero is allowed; the displayed convention therefore requires confirmation.
\subsection{Short English statement}
Does Minkowski addition of a ball preserve circle-freeness for every convex body in \(\R^3\)? The supplied record leaves circle-freeness undefined, so a complete precise statement requires source clarification. The record labels the problem open; this status has not been independently verified.
\subsection{Sources}
\begin{itemize}
\item Supplied local UnsolvedMath registry record, permanent ID \(30001883\), title \emph{Circle Free Convex Bodies Under Minkowski Addition}, source problem code \texttt{catalogue entry}, in the Oberwolfach Reports Open Problems collection. The record gives proposed year 2011 and names no proposer.
\item Supplied source citation: \emph{Discrete Geometry (2012)}, Oberwolfach Reports, DOI 10.4171/owr/2011/44, pp. 2459--2548. \url{https://doi.org/10.4171/owr/2011/44}. This is also the supplied literature source labeled ``original OWR report.'' The background describes Oberwolfach Reports as documenting mathematical workshops and their open-problem discussions.
\item Supplied literature source labeled ``circle-free convex bodies'': \url{https://www.fields.utoronto.ca/programs/scientific/11-12/discretegeom/talks/hiroshi.pdf}. No author or further bibliographic details are supplied for this link.
\item Literature assessment dated exactly \texttt{2026-08-21}: the supplied curation labels the problem open and states that known results for planar sets and special bodies do not settle the question for arbitrary circle-free convex bodies. This is supplied dated literature triage derived from the OWR curation, not an independently verified assessment; time-sensitive status and cited claims require independent checking.
\item The supplied statement status is \texttt{corrected\_verified}; its accompanying verification describes only enclosing TeX notation in mathematical delimiters without changing mathematical content. It supplies no definition of circle-freeness.
\end{itemize}
\end{document}
